\documentclass[12pt]{article}
\usepackage[a4paper,margin=1in]{geometry}
\usepackage[utf8]{inputenc}
\usepackage[T1]{fontenc}
\usepackage{amsmath,amssymb}
\usepackage{textcomp}
\usepackage{parskip}
\begin{document}
\begin{center}
{\Large \textbf{2006 STEP II}}
\end{center}

\subsection*{Page 1}
\begin{flushleft}
STEP II, 2006 2\par
\bigskip\textbf{Section A: Pure Mathematics}\par
\bigskip\textbf{Question 1.} The sequence of real numbers u1, u2, u3, . . . is defined by\par
\[u^1 = 2 , and u n+1 = k for n > 1, ( )\]
- u ∗\par
n\par
where k is a constant.\par
\textbf{(i)} Determine the values of k for which the sequence ( ) is:\par
∗\par
(a) constant;\par
(b) periodic with period 2;\par
(c) periodic with period 4.\par
\textbf{(ii)} In the case k = 37, show that u > 2 for all n. Given that in this case the sequence ( )\par
n\par
∗\par
converges to a limit `, find the value of `.\par
2 Using the series\par
\[x^2 x^3 x^4\]
\[ex = 1 + x + + + + ,\]
2! 3! 4! \cdots\par
\[show that e > .\]
\[Show that n! > 2n for n > 4 and hence show that e < 67 .\]
Show that the curve with equation\par
\[y = 3e\]
2x\par
+ 14 ln(\par
\[x) , x <\]
3 - 3\par
\[has a minimum turning point between x = and x = 1 and give a sketch to show the shape\]
of the curve.\par
\end{flushleft}
\subsection*{Page 2}
\begin{flushleft}
STEP II, 2006 3\par
\bigskip\textbf{Question 3.} (i) Show that 5 + \sqrt24 4 + 1 is an integer.\par
\[5 + \sqrt\{24\}\]
Show also that\par
1 2\par
\[0.1 < < < 0.11 .\]
\[5 + \sqrt\{24\} 19\]
\[Hence determine, with clear reasoning, the value of 5 + \sqrt\{24\} correct to four decimal\]
places.\par
k\par
\textbf{(ii)} If N is an integer greater than 1, show that N + \sqrtN 2 1 , where k is a positive\par
\gamma - \delta\par
1 -k\par
integer, differs from the integer nearest to it by less than 2N .\par
- 2\par
\bigskip\textbf{Question 4.} By making the substitution x = \pi t , show that\par
-\par
\[\pi \pi\]
\[xf(sin x)dx = \pi f(sin x)dx ,\]
Z 2 Z\par
0 0\par
where f(sin x) is a given function of sin x.\par
Evaluate the following integrals:\par
\[\pi x sin x\]
\textbf{(i)} dx ;\par
Z 0 3 + sin 2 x\par
\[2\pi\]
x sin x\par
\textbf{(ii)} dx ;\par
Z 0 3 + sin 2 x\par
\[\pi x sin 2x\]
\textbf{(iii)} dx .\par
Z 0 3 \{ \{+ sin 2 x \{ \{\par
\bigskip\textbf{Question 5.} The notation bx denotes the greatest integer less than or equal to the real number x. Thus,\par
\[for example, \pi = 3 , 18 = 18 and 4.2 = 5 .\]
b c b c b- c -\par
\textbf{(i)} Two curves are given by y = x2 + 3x 1 and y = x2 + 3 x 1 . Sketch the curves,\par
- b c -\par
for 1 6 x 6 3 , on the same axes.\par
Find the area between the two curves for 1 6 x 6 n, where n is a positive integer.\par
\textbf{(ii)} Two curves are given by y = x2 + 3x 1 and y = x 2 + 3 x 1 . Sketch the curves,\par
- b c b c -\par
for 1 6 x 6 3 , on the same axes.\par
Show that the area between the two curves for 1 6 x 6 n, where n is a positive integer, is\par
(n 1)(3n + 11) .\par
6 -\par
\end{flushleft}
\subsection*{Page 3}
\begin{flushleft}
STEP II, 2006 4\par
\bigskip\textbf{Question 6.} By considering a suitable scalar product, prove that\par
(ax + by + cz) 2 6 (a 2 + b 2 + c 2 )(x 2 + y 2 + z 2 )\par
for any real numbers a, b, c, x, y and z. Deduce a necessary and sufficient condition on a, b,\par
c, x, y and z for the following equation to hold:\par
2 2 2 2 2 2 2\par
\[(ax + by + cz) = (a + b + c )(x + y + z ) .\]
\textbf{(i)} Show that (x + 2y + 2z)2 6 9(x2 + y2 + z2) for all real numbers x, y and z.\par
\textbf{(ii)} Find real numbers p, q and r that satisfy both\par
2 2 2\par
\[p + 4q + 9r = 729 and 8p + 8q + 3r = 243 .\]
\[x^2 y^2\]
\bigskip\textbf{Question 7.} An ellipse has equation + = 1. Show that the equation of the tangent at the point\par
\[a^2 b^2\]
(a cos α, b sin α) is\par
b cot α\par
\[y = x + b cosec α .\]
- a\par
The point A has coordinates ( a, b), where a and b are positive. The point E has coordinates\par
- -\par
( a, 0) and the point P has coordinates (a, kb), where 0 < k < 1. The line through E parallel\par
-\par
\[to AP meets the line y = b at the point Q. Show that the line P Q is tangent to the above\]
\[ellipse at the point given by tan(α/2) = k.\]
\[Determine by means of sketches, or otherwise, whether this result holds also for k = 0 and\]
\[k = 1.\]
\bigskip\textbf{Question 8.} Show that the line through the points with position vectors x and y has equation\par
\[r = (1 α)x + αy ,\]
-\par
where α is a scalar parameter.\par
\[The sides OA and CB of a trapezium OABC are parallel, and OA > CB. The point E on OA\]
\[is such that OE : EA = 1 : 2, and F is the midpoint of CB. The point D is the intersection\]
of OC produced and AB produced; the point G is the intersection of OB and EF ; and the\par
point H is the intersection of DG produced and OA. Let a and c be the position vectors of\par
the points A and C, respectively, with respect to the origin O.\par
\textbf{(i)} Show that B has position vector λa + c for some scalar parameter λ.\par
\textbf{(ii)} Find, in terms of a, c and λ only, the position vectors of D, E, F , G and H. Determine\par
the ratio OH : HA.\par
\end{flushleft}
\subsection*{Page 4}
\begin{flushleft}
STEP II, 2006 5\par
\bigskip\textbf{Section B: Mechanics}\par
\bigskip\textbf{Question 9.} A painter of weight kW uses a ladder to reach the guttering on the outside wall of a house.\par
The wall is vertical and the ground is horizontal. The ladder is modelled as a uniform rod of\par
weight W and length 6a.\par
The ladder is not long enough, so the painter stands the ladder on a uniform table. The table\par
has weight 2W and a square top of side a with a leg of length a at each corner. The foot of\par
the ladder is at the centre of the table top and the ladder is inclined at an angle arctan 2 to\par
the horizontal. The edge of the table nearest the wall is parallel to the wall.\par
The coefficient of friction between the foot of the ladder and the table top is . The contact\par
between the ladder and the wall is sufficiently smooth for the effects of friction to be ignored.\par
\textbf{(i)} Show that, if the legs of the table are fixed to the ground, the ladder does not slip on\par
the table however high the painter stands on the ladder.\par
\textbf{(ii)} It is given that k = 9 and that the coefficient of friction between each table leg and the\par
ground is . If the legs of the table are not fixed to the ground, so that the table can\par
tilt or slip, determine which occurs first when the painter slowly climbs the ladder.\par
[Note: arctan 2 is another notation for tan\par
-1\par
2.]\par
\bigskip\textbf{Question 10.} Three particles, A, B and C, of masses m, km and 3m respectively, are initially at rest lying\par
in a straight line on a smooth horizontal surface. Then A is projected towards B at speed u.\par
After the collision, B collides with C. The coefficient of restitution between A and B is and\par
the coefficient of restitution between B and C is .\par
\textbf{(i)} Find the range of values of k for which A and B collide for a second time.\par
\textbf{(ii)} Given that k = 1 and that B and C are initially a distance d apart, show that the time\par
60d\par
that elapses between the two collisions of A and B is .\par
13u\par
\end{flushleft}
\subsection*{Page 5}
\begin{flushleft}
STEP II, 2006 6\par
\bigskip\textbf{Question 11.} A projectile of unit mass is fired in a northerly direction from a point on a horizontal plain at\par
speed u and an angle θ above the horizontal. It lands at a point A on the plain. In flight, the\par
projectile experiences two forces: gravity, of magnitude g; and a horizontal force of constant\par
magnitude f due to a wind blowing from North to South. Derive an expression, in terms of\par
u, g, f and θ for the distance OA.\par
\textbf{(i)} Determine the angle α such that, for all θ > α, the wind starts to blow the projectile\par
back towards O before it lands at A.\par
\textbf{(ii)} An identical projectile, which experiences the same forces, is fired from O in a northerly\par
\textasciicircum{}\circ\par
direction at speed u and angle 45 above the horizontal and lands at a point B on the\par
plain. Given that θ is chosen to maximise OA, show that\par
OB g f\par
\[= - .\]
\[OA g^2 + f 2 f\]
-\par
p\par
\[Describe carefully the motion of the second projectile when f = g.\]
\end{flushleft}
\subsection*{Page 6}
\begin{flushleft}
STEP II, 2006 7\par
\bigskip\textbf{Section C: Probability and Statistics}\par
\bigskip\textbf{Question 12.} A cricket team has only three bowlers, Arthur, Betty and Cuba, each of whom bowls 30 balls\par
in any match. Past performance reveals that, on average, Arthur takes one wicket for every 36\par
balls bowled, Betty takes one wicket for every 25 balls bowled, and Cuba takes one wicket for\par
every 41 balls bowled.\par
\textbf{(i)} In one match, the team took exactly one wicket, but the name of the bowler was not\par
recorded. Using a binomial model, find the probability that Arthur was the bowler.\par
\textbf{(ii)} Show that the average number of wickets taken by the team in a match is approxi-\par
mately 3. Give with brief justification a suitable model for the number of wickets taken\par
by the team in a match and show that the probability of the team taking at least five\par
wickets in a given match is approximately .\par
\[[You may use the approximation e^3 = 20.]\]
\bigskip\textbf{Question 13.} I know that ice-creams come in n different sizes, but I don't know what the sizes are. I am\par
offered one of each in succession, in random order. I am certainly going to choose one - the\par
bigger the better - but I am not allowed more than one. My strategy is to reject the first\par
ice-cream I am offered and choose the first one thereafter that is bigger than the first one I\par
was offered; if the first ice-cream offered is in fact the biggest one, then I have to put up with\par
the last one, however small.\par
\[Let P (k) be the probability that I choose the kth biggest ice-cream, where k = 1 is the\]
n\par
\[biggest and k = n is the smallest.\]
\textbf{(i)} Show that P4(1) =\par
\[and find P^4(2), P^4(3) and P^4(4).\]
\textbf{(ii)} Find an expression for P (1).\par
n\par
\end{flushleft}
\subsection*{Page 7}
\begin{flushleft}
STEP II, 2006 8\par
\bigskip\textbf{Question 14.} Sketch the graph of y = for x > 0, x = 1. You may assume that x ln x 0 as x 0.\par
x ln x 6 \to \to\par
The continuous random variable X has probability density function\par
λ\par
for a 6 x 6 b ,\par
\[f(x) =  x ln x\]
\par
\par
0 otherwise ,\par
\par
\par
where a, b and λ are suitably chosen constants.\par
\textbf{(i)} In the case a = 1/4 and b = 1/2, find λ .\par
\textbf{(ii)} In the case λ = 1 and a > 1, show that b = ae.\par
\textbf{(iii)} In the case λ = 1 and a = e, show that P(e3/2 6 X 6 e2) 31 .\par
\approx 108\par
\textbf{(iv)} In the case λ = 1 and a = e1/2, find P(e3/2 6 X 6 e2) .\par
\end{flushleft}
\end{document}
